\section{Virtualização De Dispositivos}

Dispositivos externos em sistemas RISC-V são, em sua maioria, acessados por meio de \gls{mmio}, no qual registradores de controle e dados dos dispositivos são mapeados em endereços do
espaço de memória. Dessa forma, uma operação de leitura ou escrita em um determinado endereço pode representar uma operação sobre um dispositivo, em vez de um acesso à memória convencional.
A virtualização desses dispositivos pode, portanto, ser realizada interceptando as operações de memória destinadas às regiões correspondentes aos dispositivos virtuais e emulando seu comportamento.

Cada máquina virtual possui uma região de memória própria. Os acessos do convidado a essa região são permitidos pela \gls{pmp}, enquanto as demais regiões de memória permanecem protegidas. Assim,
quando o sistema convidado tenta acessar um endereço que não pertence à sua região de memória, o processador gera uma exceção de acesso, como \textit{Load Access Fault} ou \textit{Store Access Fault},
e transfere o controle para o hipervisor.

O hipervisor utiliza essas exceções para distinguir acessos inválidos à memória de acessos destinados a dispositivos virtuais. Quando o endereço acessado corresponde à região de \gls{mmio} de
um dispositivo virtual, o hipervisor decodifica a instrução que originou a exceção e identifica o tipo da operação, o endereço acessado, seu tamanho e, no caso de uma escrita,
o valor que deveria ser armazenado.

No caso de uma operação de leitura, o hipervisor solicita ao dispositivo virtual correspondente que realize a operação. O dispositivo virtual processa a requisição de
acordo com seu estado interno e retorna o valor correspondente. O hipervisor então escreve esse valor no registrador de destino do contexto do convidado e ajusta o
contador de programa para que a instrução que causou a exceção não seja executada novamente.

\begin{figure}[H]
\centering
\includegraphics[width=0.8\textwidth]{images/LoadAccessFault.pdf}
\caption{Emulação de Dispositivos via \textit{Load Access Fault}}\label{fig:emulate-load-access-fault}
\end{figure}

Para operações de escrita, o processo é semelhante. O hipervisor identifica o valor que seria escrito pelo convidado e o encaminha ao dispositivo virtual correspondente.
O dispositivo pode atualizar seu estado interno, produzir efeitos colaterais ou encaminhar a operação para o mediador de hardware correspondente.

\begin{figure}[H]
\centering
\includegraphics[width=0.65\textwidth]{images/StoreAccessFault.pdf}
\caption{Emulação de Dispositivos via \textit{Store Access Fault}}\label{fig:emulate-store-access-fault}
\end{figure}

Caso o endereço não corresponda a nenhum dispositivo virtual ou a operação solicitada não seja suportada pelo dispositivo, o acesso não pode ser emulado.
Nesse caso, o hipervisor pode tratar a operação como um acesso inválido e encaminhar uma exceção ao sistema convidado ou, dependendo da política adotada, encerrar a \gls{vm}.
